\documentclass[11pt]{article}
\usepackage[utf8]{inputenc}
\usepackage[T1]{fontenc}
\usepackage[italian]{babel}
\usepackage{geometry}
\geometry{a4paper, margin=1in}
\usepackage{xcolor}
\usepackage{listings}
\usepackage{hyperref}

\lstset{
    language=C,
    basicstyle=\ttfamily\small,
    keywordstyle=\color{blue},
    stringstyle=\color{red},
    commentstyle=\color{green!60!black},
    numbers=none,
    frame=single,
    breaklines=true
}

\usepackage{tcolorbox}
\tcbuselibrary{breakable}
\begin{document}

\begin{tcolorbox}[colback=blue!5!white,colframe=blue!75!black,title=\textbf{DOMANDA}, breakable]
Temi d'esame (per ognuno ragionare sulla scelta del protocollo di livello trasporto)
\begin{itemize}
    \item Implementare un servizio di calcolo remoto: il client spedisce al server in un unico messaggio due numeri interi e il tipo di operatore (somma, sottrazione, moltiplicazione, divisione) e il server, dopo aver svolto il calcolo, restituisce il risultato al client che così termina la sua esecuzione.
\end{itemize}
\end{tcolorbox}

\begin{tcolorbox}[colback=green!5!white,colframe=green!75!black,title=\textbf{RISPOSTA}, breakable]
\textbf{Piano di Implementazione:} Servizio di Calcolo Remoto (es17)

\textbf{1. Scelta del Protocollo di Trasporto}\\
La traccia chiede di ragionare sulla scelta tra \texttt{UDP} e \texttt{TCP} per un servizio in cui il client invia un unico messaggio con i dati, riceve la risposta e termina.

\textbf{Scelta ideale:} \texttt{UDP}

\begin{itemize}
    \item \textbf{Motivazione:} Il modello di comunicazione richiesto è strettamente di tipo "Richiesta-Risposta" singola. Il payload (due interi e un carattere per l'operatore) è minuscolo e rientra comodamente in un unico datagramma \texttt{UDP}.
    \item \textbf{Efficienza:} Usare \texttt{TCP} richiederebbe un notevole overhead per stabilire la connessione (3-way handshake) e abbatterla, scambiando una decina di pacchetti di controllo solo per inviare due messaggi utili. Con \texttt{UDP}, i pacchetti in gioco sono solo due (uno di andata, uno di ritorno).
    \item \textbf{Risoluzione di inaffidabilità:} Trattandosi di una singola query, se il pacchetto va perso il client potrebbe semplicemente riprovare (con un timeout applicativo), in maniera del tutto analoga a come funziona il protocollo DNS.
\end{itemize}
Pertanto, l'implementazione proposta utilizzerà \texttt{UDP}.

\textbf{2. Struttura dei Dati}\\
Definiremo una struct per il messaggio di richiesta:

\begin{lstlisting}
typedef struct {
    int operando1;
    int operando2;
    char operatore; // '+', '-', '*', '/'
} RichiestaCalcolo;
\end{lstlisting}

\textbf{3. Implementazione del Client (\texttt{clientUDP.c})}
\begin{itemize}
    \item Crea un'interfaccia \texttt{UDP} (es. porta 20000).
    \item Chiede in input all'utente i due numeri e l'operatore.
    \item Popola la struct \texttt{RichiestaCalcolo}.
    \item Spedisce la struct al server (\texttt{UDPSend}).
    \item Attende il risultato in ricezione (\texttt{UDPReceive}).
    \item Stampa il risultato e termina.
\end{itemize}

\textbf{4. Implementazione del Server (\texttt{serverUDP.c})}
\begin{itemize}
    \item Crea un'interfaccia \texttt{UDP} in ascolto (es. porta 35000).
    \item Entra in un ciclo infinito di ricezione \texttt{UDPReceive} in attesa di richieste.
    \item Quando riceve una \texttt{RichiestaCalcolo}, valuta l'operatore (\texttt{switch} statement).
    \item Calcola il risultato (float o intero, gestiremo le divisioni o daremo un intero).
    \item Risponde al client inviando il risultato.
\end{itemize}
\end{tcolorbox}

\end{document}
